###**1.**引言

自古以来，人类对于自身大脑的探索从未停止。大脑，这个复杂而精密的器官，不仅赋予了我们思考、感知和学习的能力，更让我们在漫长的进化过程中脱颖而出，成为地球上最智慧的生物。然而，尽管科学家们已经取得了许多关于大脑结构和功能的重大发现，但大脑的工作机制仍然充满了未知和神秘。

在大脑的探索之旅中，一个尤为引人注目的领域是神经网络。神经网络，特别是人脑中的神经网络，是一个由数以亿计的神经元通过突触相互连接而成的庞大网络。这些神经元和突触不仅负责传递和处理信息，还通过复杂的交互作用，实现了我们的各种感知、思考和行为。

受到人脑神经网络的启发，科学家们开始尝试构建人工神经网络（ArtificialNeuralNetworks,ANNs），以模拟大脑的信息处理过程。

人工神经网络 4大理论支柱为 “阈值逻辑”、“Hebb 学习率”、“梯度下降”、“反向传播”，前2个理论解决了单个神经元层面的建模问题，后2个理论则解决了多层神经网络训练问题。

2006年 Hinton 首次实现了5 层神经网络的训练，之后行业迎来爆发式发展，不断验证了该技术。

深度学习具备坚实的数学理论基础支撑。人脑绝大多数活动本质上都是广义计算问题，因此人脑其实是一个复杂的函数，深度学习就是去找到这个函数，万能近似定理则从数学上证明了一定条件下深度神经网络模型能够模拟任意的函数。
用于深度学习的算力以6年30 万倍的速度增加，算力是核心瓶颈也是未来提升的关键。从进化角度看，人的智能是一个随着神经元数量提升，从量变到质变的过程，这个量变过程对应人工智能中模型所用算力的提升过程。目前能够实现的 AI 模型中不论从神经元个数还是连接数量看，与真实人类还有较大的差距，未来随着现有芯片技术的不断推进，和突破冯诺依曼架构的类脑芯片等新技术的发展，算力的持续突破将会不断释放人工智能技术的潜力。


**面对大数据的人类大脑**

这些算法似乎都离日常生活非常遥远。你可能会说你能够一一记住自己接收的所有数据，但这就错得离谱了。

人们可能从未察觉，每时每刻充斥我们大脑皮层的数据量大得可怕，有如CERN需要处理的数据量。我们的视觉、听觉、嗅觉、触觉、温度感觉，以及其他数不清的感觉，每秒都会向我们传输大约十亿字节(1GB)的数据[4]。如果我们把数十年间收集到的数据累计起来的话，我们会发现感官接收到的数据总量需要用百亿亿字节(EB)作为计量单位!这可以说是大得可怕。

与现代信息技术面对的情况类似，我们的大脑也无法储存自身接收到的所有数据。它也不想这样做，因为其中绝大部分数据没有意义，它不得不忘记绝大部分数据。

在视觉数据的处理中，这种状况尤其突出。据神经科学家马库斯·赖希勒所言[5]，我们的视皮层拥有惊人的能力，可以将十亿字节的视觉数据转换为几千字节的有用数据，这种处理在每一秒都在进行，而且只需要消耗极少的能量[^footnote视觉]。
[^footnote视觉]:赖希勒估计，我们的视网膜每秒会收集到大约10^10比特的数据，但到达初级视皮层的第四层的数据只有大约10^4比特。


在更普遍的意义上，大脑会优先保留数据的“大体概念”，而不是具体细节。一般来说，我敢打赌你不能背出这本书前15章中的任意一句话。但在我的期望中，你应该记住了这些章节讨论的是贝叶斯公式、它的逻辑基础、它在归纳推理问题上的应用、它的历史、所罗门诺夫妖及其在博弈论中的应用，甚至还有它与演化理论之间的联系。即使你可能没有记住读到的任何一个句子，但我仍然希望你记住了读到的文字中各种基础概念的抽象表达。

对于接收到的信息只保留压缩过后的表示，这种能力并非弱点，而正是大脑的强大之处。我们一般都能够回忆起那些重要的事物，而且完全忘却那些对我们来说无足轻重的东西。

**用贝叶斯帮助记忆**
我们之前看到，人类大脑令人叹服的地方之一，就是它将极大量原始数据压缩为寥寥几个想法的能力。受此启发，人工智能研究者发明了
自编码器这一架构。
自编码器的任务，就是对大量信息进行精简，这些信息可以是高分辨率的图像或者整部电影，等等，最终只留下内容的精髓。换句话说，
这些神经网络尝试做到的，正是我们在语文课上被要求做的一种练习:撰写摘要。为了测试摘要的质量，人们也要求自编码器对自己生
成的摘要进行解压缩，设想与这一摘要相关的高分辨率图像或其他原始数据是什么样子的。
要做到这一点，关键之一就是利用贝叶斯方法。对数据重组，其实就是确定什么数据会得出眼前这个摘要。这正是贝叶斯公式的典型应
用!我们需要确定这个摘要的可能原因，因此

<!-- ![截屏2024-09-0418.50.53](/Users/Min369/Downloads/同步空间/Obsidian/书/截屏2024-09-0418.50.53.png) -->

另外，我们在下一章也会看到人工智能研究者如何利用这一公式，引入不常见的摘要，向机器赋予了创造性。

**短期记忆与长期记忆**

金特和朱莉娅·肖研究的是长期记忆。其中，信息被储存在大脑皮层这一神经网络的拓扑结构之中。问题在于访问这些数据可能很困难，需要测试信号在网络中传播的各种可能方式。回忆一首有名的歌曲的歌词一般来说就是这样的情况。在刚开始回忆时，我们通常想起几个词之后就想不起来了，就好像大脑中电信号的流动停滞了，或者流向了错误的方向一样。然而在重复这种流动之后，我们最终仍能找到正确的道路，回忆起整首歌的歌词。

人们在大型数据库中搜寻信息的时候也会碰到同样的问题。给谷歌和其他公司带来巨大财富的事物之一，就是为了加速信息搜索而整合互
联网数据的方法。但这并不只是信息整合的问题，人们也经常要在储存容量和存储速度之间对储存信息的媒介进行取舍。

**递归神经网络**

无论如何，人工神经网络的研究选择了第三条道路:神经信号的环路传播。包含这种环路的神经网络架构又叫**递归神经网络**。在处理当前数据时，这种网络能够利用前一时刻的部分信息。对于处理本质上有时序特性的数据，比如书中的文本或者讲话中的声音，递归神经网络已经成为最尖端的技术。

递归神经网络用到的技巧也包含在一类更广泛的拥有内部状态的算法之中。值得一提的是，这一类算法包括之前谈到的卡尔曼滤波器和隐马尔可夫模型。内部状态的任务就是以精练的方式提取出前一时刻的信息，用以更好地理解前一个被分析的数据。然而另一个问题就是，与数据读取相比，内部状态会以什么样的速度变化。
我有幸在自己学习数学和指导学生学习的过程中都观察到了一点东西。第一次读一份文献时(比如说你正在看的这本书)，最好不要停留在细节上，以免偏离阅读的主线。因为如果读得太慢，换种说法，就是对于短期记忆这一内部状态修改得太多的话，会使思绪偏离文献
想让我们思考的东西。因此，有时候先迅速浏览一遍，再多花点时间细细阅读，最后再速读一遍，这样可能效果更好。每一次阅读都会带来新的东西，因为每一次阅读都联系着短期记忆这一内部状态的不同变化动态。
到这里为止，我们讨论的还是对数据的线性阅读，但我们阅读或者聆听的方式并非完全线性的。的确，当某个句子非常晦涩的时候，我们一般会重读几次，然后才继续线性阅读。同样，有时候如果别人一句话还没说完，我们就无法理解。在梵语、印地语和日语中，动词位于句末，所以这种情况尤其显著。同样道理，在数学论文中，计算过程之后通常会有对计算的解释。还有些笑话、广告和电影，正是结尾决定了它们的意义。

为了同时利用过去和未来的数据理解现在，人们提出了另一种神经网络架构，那就是所谓的**双向递归神经网络**。为了将未来纳入其中，这些网络必须延长得出回应的时间。我们可以打赌，人类大脑中也有类似的结构。这也能解释我们在经过一段延迟之后才理解笑话的能力，甚至有时候这种延迟本身就引人发笑。

最后，人工智能领域的一项最新进展就是引入了迫使遗忘的神经元件。那就是所谓的LSTM架构，意思是“长短期记忆”(long-short-termmemory)。除了递归神经网络中常见的那些神经元环路之外，LSTM还拥有另一个额外的环路，它激活的时候会强制让所有神经元环路中的信号消失。这也许就解释了为什么我们在讨论中被打断之后，有时候很难回想起之前讨论的话题。

目前，在处理本质上有时序特性的数据方面，比如语音识别、自然语言分析，等等，LSTM及其变体似乎就是最尖端的人工智能方法。但我承认自己不理解它为何如此成功，我也承认自己没有足够的专业水平来预测这些领域未来的前沿进展。


###2.背景



1.####第一阶段

1943年，美国神经生理学家沃伦·麦卡洛克和数学家沃尔特·皮茨对生物神经元进行建模，首次提出了一种形式神经元模型。这个神经元模型通过电阻等元件构建的物理网络得以实现，被称为M-P模型。1958年，罗森布拉特又提出来了感知器，这意味着经过训练后，计算机能够确定神经元的连接权重。就这样，神经网络的研究迎来了第一次高潮。然而在1969年，明斯基等人指出感知器无法解决线性不可分问题，使得神经网络网络的研究陷入了第一次低谷。



2.####第二阶段

1986年，Rumelhart(鲁梅尔哈特)等人提出了误差反向传播算法通过设置多层感知器，解决了线性不可分问题。从反向传播算法的发表开始，神经网络迎来第二次高潮。1989年，YannLecun等研究人员实现了一个7层的卷积神经网络LeNet-5，识别手写数字的精度达到了98%。在这段时间内，一大批学者和研究人员围绕Hopfield提出的人工神经网络方法展开了进一步研究，形成了20世纪80年代中期以来人工神经网络的研究热潮。可是好景不长，在20世纪末，支持向量机(SVM)和其他机器学习算法的流行程度逐渐高涨，更重要的是其他机器学习算法在实现效率上远超神经网络。随着PageRank等互联网新兴的算法不断被提出，学者们的研究重心也开始转移到其他机器学习算法上，神经网络的研究热潮又迎来了第二次低谷。



3.####第三阶段

2006年，Hinton(辛顿)等人利用限制玻尔兹曼机对神经网络的连续层进行建模，使用逐层预训练的方法抽取模型数据中的高维特征，后来又提出了深度信念网络。这一技巧随后被众多学者推广到了许多不同的神经网络架构上，大大地提高了模型在测试集上的泛化效果。同时，随着硬件计算能力和算法的提升，网络隐层可以不断增加，于是以Hinton为代表的研究人员重新将人工神经网络定义为深度学习人工神经网络定义为深度学习，并进行推广。自201l年起，神经网络就在语音识别和图像识别基准测试中获得了压倒性优势，由此迎来了它的第三次高潮。





![img](file:////Users/Min369/Library/Group%20Containers/UBF8T346G9.Office/TemporaryItems/msohtmlclip/clip_image002.jpg)









###3.结构

![R-C](file:////Users/Min369/Library/Group%20Containers/UBF8T346G9.Office/TemporaryItems/msohtmlclip/clip_image004.jpg)

三个层次：输入层、隐藏层和输出层

输入层是神经网络的基层，它负责接收外界的输入数据。这些数据可以是图像、声音、文本等各种形式，经过预处理后转换成神经网络能够接受的形式。
隐藏层是神经网络的核心部分，它由多个神经元组成，负责处理输入数据并传递到下一层。隐藏层的设计和数量需要根据具体的任务来确定，它们的作用是对输入数据进行分类或预测。
输出层是神经网络的最高层，它负责将隐藏层处理后的结果转换成实际输出



圆形节点与人工神经元：

在人工神经网络中，每个圆形节点代表一个人工神经元。这些神经元通过特定的连接方式相互交互，模拟生物神经网络的工作原理。



连接与信号传递：

箭头表示从一个神经元的输出到另一个神经元的输入的连接。通过这些连接，信号可以在网络中传递，从一个人工神经元传递到另一个。



权重与激活函数：

每个节点都代表一种特定的输出函数，称为激活函数。每两个节点间的连接都有一个与之相关的权重值，表示前一个神经元对后一个神经元的影响程度。



网络输出：

网络的输出会根据网络的连接方式、权重值以及激励函数的不同而变化。通过调整这些参数，人工神经网络能够学习和适应不同的输入模式，产生预期的输出结果。



###4.如何训练

**4.1****简介**

反向传播算法（Backpropagation）是目前用来训练人工神经网络（ArtificialNeuralNetwork，ANN）的最常用且最有效的算法。其主要思想是：

（1）将训练集数据输入到ANN的输入层，经过隐藏层，最后达到输出层并输出结果，这是ANN的前向传播过程；

（2）由于ANN的输出结果与实际结果有误差，则计算估计值与实际值之间的误差，并将该误差从输出层向隐藏层反向传播，直至传播到输入层；

（3）在反向传播的过程中，根据误差调整各种参数的值；不断迭代上述过程，直至收敛。

**4.2****数学推导**

**1.****变量定义**

先定义一些变量：

​![IMG_256](file:////Users/Min369/Library/Group%20Containers/UBF8T346G9.Office/TemporaryItems/msohtmlclip/clip_image005.jpg)表示第l-1层的第k个神经元连接到第l层的第j个神经元的权重；

​![IMG_256](file:////Users/Min369/Library/Group%20Containers/UBF8T346G9.Office/TemporaryItems/msohtmlclip/clip_image006.jpg)表示第l层的第j个神经元的偏置；

![IMG_256](file:////Users/Min369/Library/Group%20Containers/UBF8T346G9.Office/TemporaryItems/msohtmlclip/clip_image007.jpg)表示第l层的第j个神经元的输入，即：

​![IMG_256](file:////Users/Min369/Library/Group%20Containers/UBF8T346G9.Office/TemporaryItems/msohtmlclip/clip_image008.jpg)

![IMG_256](file:////Users/Min369/Library/Group%20Containers/UBF8T346G9.Office/TemporaryItems/msohtmlclip/clip_image009.jpg)表示第l层的第j个神经元的输出，即：

​![IMG_256](file:////Users/Min369/Library/Group%20Containers/UBF8T346G9.Office/TemporaryItems/msohtmlclip/clip_image010.jpg)



​![IMG_256](file:////Users/Min369/Library/Group%20Containers/UBF8T346G9.Office/TemporaryItems/msohtmlclip/clip_image011.jpg)其中表示激活函数。



2.**代价函数**

​代价函数被用来计算ANN输出值与实际值之间的误差。常用的代价函数是二次代价函数（Quadraticcostfunction）：

​![IMG_256](file:////Users/Min369/Library/Group%20Containers/UBF8T346G9.Office/TemporaryItems/msohtmlclip/clip_image012.jpg)





​其中，x表示输入的样本，y表示实际的分类，![IMG_256](file:////Users/Min369/Library/Group%20Containers/UBF8T346G9.Office/TemporaryItems/msohtmlclip/clip_image013.jpg)表示预测的输出，L表示神经网络的最大层数。



3.**公式及其推导**

本节将介绍反向传播算法用到的4个公式，并进行推导。如果不想了解公式推导过程，请直接看第4节的算法步骤。

![IMG_256](file:////Users/Min369/Library/Group%20Containers/UBF8T346G9.Office/TemporaryItems/msohtmlclip/clip_image014.jpg)首先，将第层第个神经元中产生的错误（即实际值与预测值之间的误差）定义为：





公式1（计算最后一层神经网络产生的错误）：



​![IMG_256](file:////Users/Min369/Library/Group%20Containers/UBF8T346G9.Office/TemporaryItems/msohtmlclip/clip_image016.jpg)

其中，![IMG_256](file:////Users/Min369/Library/Group%20Containers/UBF8T346G9.Office/TemporaryItems/msohtmlclip/clip_image017.jpg)表示Hadamard乘积，用于矩阵或向量之间点对点的乘法运算。公式1的推导过程如下：



![IMG_256](file:////Users/Min369/Library/Group%20Containers/UBF8T346G9.Office/TemporaryItems/msohtmlclip/clip_image019.jpg)

首先看上面证明的第一行，我们知道损失函数的形式，输出层的每一个输出都能够看作是损失函数的一个变量。

要是不理解的话，假设我们这里的损失函数为之前提到过的二次函数

在神经网络里面可以写为上面的形式，因为一个网络的输出就是最后一层的输出。这下能不能够明显的理解上面的那句话了呢？

那么证明第一行就是通过链式法则得到的。

然后只有![IMG_256](file:////Users/Min369/Library/Group%20Containers/UBF8T346G9.Office/TemporaryItems/msohtmlclip/clip_image021.jpg)和![IMG_256](file:////Users/Min369/Library/Group%20Containers/UBF8T346G9.Office/TemporaryItems/msohtmlclip/clip_image023.jpg)是有依赖关系的，其他的没有依赖关系的求导为0不见了，因此，只剩下第二行。第二行出来了，又因为![IMG_256](file:////Users/Min369/Library/Group%20Containers/UBF8T346G9.Office/TemporaryItems/msohtmlclip/clip_image025.jpg)，那么后面的也就都容易推出来了.





公式2（由后往前，计算每一层神经网络产生的错误）：

​![IMG_256](file:////Users/Min369/Library/Group%20Containers/UBF8T346G9.Office/TemporaryItems/msohtmlclip/clip_image027.jpg)

推导过程：

​![IMG_256](file:////Users/Min369/Library/Group%20Containers/UBF8T346G9.Office/TemporaryItems/msohtmlclip/clip_image029.jpg)



通过这个公式，只要我知道l+1层的误差，那么我就能够知道l层的误差，由此，我就能够倒退直到第1层的误差。
这个式子就体现了误差的反向传播的特点。
然后公式1和公式2结合起来，我们就能够求出神经网络中任何一层，任何一个神经元上面的误差。

公式3（计算权重的梯度）：

​![IMG_256](file:////Users/Min369/Library/Group%20Containers/UBF8T346G9.Office/TemporaryItems/msohtmlclip/clip_image031.jpg)

推导过程：

​![IMG_256](file:////Users/Min369/Library/Group%20Containers/UBF8T346G9.Office/TemporaryItems/msohtmlclip/clip_image033.jpg)

这个公式告诉我们，某个神经元上面误差对于偏置的偏导要等于在这个神经元上面的误差。而我们之前已经说过了，我们要求梯度的话就需要求出这些偏导，现在通过误差间接得到了偏导。很巧妙。

公式4（计算偏置的梯度）：

​![IMG_256](file:////Users/Min369/Library/Group%20Containers/UBF8T346G9.Office/TemporaryItems/msohtmlclip/clip_image035.jpg)

推导过程：

​![IMG_256](file:////Users/Min369/Library/Group%20Containers/UBF8T346G9.Office/TemporaryItems/msohtmlclip/clip_image037.jpg)

4.反向传播算法伪代码



输入训练集



对于训练集中的每个样本x，设置输入层（Inputlayer）对应的激活值：

前向传播：![IMG_256](file:////Users/Min369/Library/Group%20Containers/UBF8T346G9.Office/TemporaryItems/msohtmlclip/clip_image039.jpg)

计算输出层产生的错误：![IMG_256](file:////Users/Min369/Library/Group%20Containers/UBF8T346G9.Office/TemporaryItems/msohtmlclip/clip_image041.jpg)

反向传播错误：![IMG_256](file:////Users/Min369/Library/Group%20Containers/UBF8T346G9.Office/TemporaryItems/msohtmlclip/clip_image043.jpg)

使用梯度下降（gradientdescent），训练参数：

![IMG_256](file:////Users/Min369/Library/Group%20Containers/UBF8T346G9.Office/TemporaryItems/msohtmlclip/clip_image045.jpg)

![IMG_256](file:////Users/Min369/Library/Group%20Containers/UBF8T346G9.Office/TemporaryItems/msohtmlclip/clip_image047.jpg)

**4.3****激活函数**

激活函数（ActivationFunction），就是在[人工神经网络](https://baike.baidu.com/item/人工神经网络/382460)的神经元上运行的[函数](https://baike.baidu.com/item/函数/301912)，负责将神经元的输入映射到输出端，旨在帮助网络学习数据中的复杂模式。

常见类型：ReLU、Sigmoid、Tanh



**Sigmoid**激活函数的数学表达式为：

​![IMG_256](file:////Users/Min369/Library/Group%20Containers/UBF8T346G9.Office/TemporaryItems/msohtmlclip/clip_image049.jpg)

导数表达式为：

![IMG_257](file:////Users/Min369/Library/Group%20Containers/UBF8T346G9.Office/TemporaryItems/msohtmlclip/clip_image051.jpg)

函数图像如下：

![IMG_258](file:////Users/Min369/Library/Group%20Containers/UBF8T346G9.Office/TemporaryItems/msohtmlclip/clip_image053.jpg)



tanh激活函数的数学表达式为：![IMG_256](file:////Users/Min369/Library/Group%20Containers/UBF8T346G9.Office/TemporaryItems/msohtmlclip/clip_image055.jpg)

函数图像如下：



![IMG_256](file:////Users/Min369/Library/Group%20Containers/UBF8T346G9.Office/TemporaryItems/msohtmlclip/clip_image057.jpg)





实际上，Tanh函数是sigmoid的变形：

与sigmoid不同的是，tanh是“零为中心”的。因此在实际应用中，tanh会比sigmoid更好一些。但是在饱和神经元的情况下，tanh还是没有解决梯度消失问题。

什么情况下适合使用Tanh？

tanh的输出间隔为1，并且整个函数以0为中心，比sigmoid函数更好；

在tanh图中，负输入将被强映射为负，而零输入被映射为接近零。

Tanh有哪些缺点？

仍然存在梯度饱和的问题

依然进行的是指数运算





**ReLU**激活函数的数学表达式为：![IMG_256](file:////Users/Min369/Library/Group%20Containers/UBF8T346G9.Office/TemporaryItems/msohtmlclip/clip_image059.jpg)

函数图像如下：

![IMG_257](file:////Users/Min369/Library/Group%20Containers/UBF8T346G9.Office/TemporaryItems/msohtmlclip/clip_image061.jpg)

什么情况下适合使用ReLU？

ReLU解决了梯度消失的问题，当输入值为正时，神经元不会饱和

由于ReLU线性、非饱和的性质，在SGD中能够快速收敛

计算复杂度低，不需要进行指数运算

ReLU有哪些缺点?

与Sigmoid一样，其输出不是以0为中心的

DeadReLU问题。当输入为负时，梯度为0。这个神经元及之后的神经元梯度永远为0，不再对任何数据有所响应，导致相应参数永远不会被更新

训练神经网络的时候，一旦学习率没有设置好，第一次更新权重的时候，输入是负值，那么这个含有ReLU的神经节点就会死亡，再也不会被激活。所以，要设置一个合适的较小的学习率，来降低这种情况的发生

**4.1****人工神经网络在各个领域中的应用**



###人工神经网络在信息领域中的应用

在处理许多问题中，信息来源既不完整，又包含假象，决策规则有时相互矛盾，有时无章可循，这给传统的信息处理方式带来了很大的困难，而神经网络却能很好的处理这些问题，并给出合理的识别与判断。

1.信息处理

现代信息处理要解决的问题是很复杂的，人工神经网络具有模仿或代替与人的思维有关的功能，可以实现自动诊断、问题求解，解决传统方法所不能或难以解决的问题。人工神经网络系统具有很高的容错性、鲁棒性及自组织性，即使连接线遭到很高程度的破坏，它仍能处在优化工作状态，这点在军事系统电子设备中得到广泛的应用。现有的智能信息系统有智能仪器、自动跟踪监测仪器系统、自动控制制导系统、自动故障诊断和报警系统等。

2.模式识别

模式识别是对表征事物或现象的各种形式的信息进行处理和分析，来对事物或现象进行描述、辨认、分类和解释的过程。该技术以贝叶斯概率论和申农的信息论为理论基础，对信息的处理过程更接近人类大脑的逻辑思维过程。现在有两种基本的模式识别方法，即统计模式识别方法和结构模式识别方法。人工神经网络是模式识别中的常用方法，近年来发展起来的人工神经网络模式的识别方法逐渐取代传统的模式识别方法。经过多年的研究和发展，模式识别已成为当前比较先进的技术，被广泛应用到文字识别、语音识别、指纹识别、遥感图像识别、人脸识别、手写体字符的识别、工业故障检测、精确制导等方面。

###人工神经网络在医学中的应用

由于人体和疾病的复杂性、不可预测性，在生物信号与信息的表现形式上、变化规律（自身变化与医学干预后变化）上，对其进行检测与信号表达，获取的数据及信息的分析、决策等诸多方面都存在非常复杂的非线性联系，适合人工神经网络的应用。目前的研究几乎涉及从基础医学到临床医学的各个方面，主要应用在生物信号的检测与自动分析，医学专家系统等。

1.医学专家系统

传统的专家系统，是把专家的经验和知识以规则的形式存储在计算机中，建立知识库，用逻辑推理的方式进行医疗诊断。但是在实际应用中，随着数据库规模的增大，将导致知识“爆炸”，在知识获取途径中也存在“瓶颈”问题，致使工作效率很低。以非线性并行处理为基础的神经网络为专家系统的研究指明了新的发展方向，解决了专家系统的以上问题，并提高了知识的推理、自组织、自学习能力，从而神经网络在医学专家系统中得到广泛的应用和发展。在麻醉与危重医学等相关领域的研究中，涉及到多生理变量的分析与预测，在临床数据中存在着一些尚未发现或无确切证据的关系与现象，信号的处理，干扰信号的自动区分检测，各种临床状况的预测等，都可以应用到人工神经网络技术。

###人工神经网络在交通领域的应用

今年来人们对神经网络在交通运输系统中的应用开始了深入的研究。交通运输问题是高度非线性的，可获得的数据通常是大量的、复杂的，用神经网络处理相关问题有它巨大的优越性。应用范围涉及到汽车驾驶员行为的模拟、参数估计、路面维护、车辆检测与分类、交通模式分析、货物运营管理、等领域并已经取得了很好的效果。

![img](file:////Users/Min369/Library/Group%20Containers/UBF8T346G9.Office/TemporaryItems/msohtmlclip/clip_image063.jpg)

**4.2****人工神经网络应用和卷积对比**

随着科技的快速发展，深度学习已经成为了人工智能领域的重要分支。在深度学习中，卷积神经网络（CNN）和人工神经网络（ANN）是两种最为常见的网络类型。它们在结构和应用上存在一些显著的区别。本节将对比分析CNN和ANN的区别，并通过实例阐述两者的应用。



CNN和ANN的区别主要体现在两个方面：结构和应用。首先，从结构上看，CNN侧重于处理图像数据，它的卷积层和池化层能够有效降低参数的数量，提高网络的性能；而ANN则没有这种结构上的限制，可以自由地设计神经网络结构来解决不同的问题。其次，从应用上看，CNN在处理图像任务时表现出色，如图像分类、目标检测等；而ANN则更加通用，可以应用于各种类型的数据和任务。
在应用案例方面，CNN和ANN各有优势。例如，在[图像识别](https://cloud.baidu.com/product/imagerecognition)任务中，CNN可以通过学习自动提取图像的特征，从而实现高效的图像分类。相比之下，ANN需要手动设计特征，因此在这方面的表现可能不如CNN。然而，在某些复杂的任务中，如[自然语言处理](https://cloud.baidu.com/product/wenxinworkshop)（NLP），ANN则可能更有优势。ANN可以通过构建深层神经网络，学习和理解语言的复杂模式和结构，从而实现高效的文本分类和语言翻译等任务。

![img](file:////Users/Min369/Library/Group%20Containers/UBF8T346G9.Office/TemporaryItems/msohtmlclip/clip_image065.jpg)





猫是什么?
2012年，谷歌因一条奇怪而又令人不安的公告登上了各大新闻头条。
据这些头条所说，谷歌的人工智能发现了“猫”这个概念[9]!很多人可能觉得这是个平平无奇的新闻，但对我来说，这是机器学习的一个里程碑式的惊人突破，可能比AlphaGo在四年之后大败李世石给人留下的印象更深刻、更出人意料。
要理解这一点，首先要说明谷歌的人工智能是一个人工神经网络，带有一些接收器，让它可以“看见”数字化图像。谷歌给这个人工神经网络展示了1000万幅图像，其中不包含上下文。然后，谷歌将这个人工神经网络放进某种相当于核磁共振成像的处理系统中，用以测量它暴露在其他图像中的时候神经元的实时激活状态。谷歌意识到其中某些神经元大体上当且仅当向其展示的图像包含一只猫时才会激活!
真正给人留下深刻印象的是，这并不是谷歌在设计这个人工智能的时候设定的目标。这个人工智能的目标是以最合理有效的方式分析、处理并解释图像中的内容。人工智能必须为它看到的那些图像建立一个模型。为此，它必须做的事情之一就是解释图像中某些像素颜色之间反复出现的相关性。比如说，当图像中的某些像素构成了眼睛的形
状，那么在它的左边或者右边不远处通常会有一份与这些像素相似的复制品。一只眼睛的图像通常伴随着另一只眼睛的图像。怎么解释眼睛很少单独出现这一点呢?
最惊人的其实是，这个问题的解释似乎大体符合我们之中大多数人会给出的答案:因为照片拍到的动物通常有两只眼睛。人、狗和猫(几乎)都拥有两只眼睛这个事实对我们来说似乎如此微不足道，人们几乎不会对它花上任何时间进行思考。我们是如此习惯于日常生活中的事物，以至于我们甚至不再寻求解释——但这其实是个迷人的生物学问题，需要对视差的数学理解!
但对我来说更迷人的是，我们竟然可以(大体上)对于什么是人、什么是猫、什么又是狗达成共识。我们甚至没有意识到这些概念既抽象又模糊!猫是什么?怎么定义猫这个概念?更重要的是，猫存在吗?
要理解“猫”这个概念的奇怪之处，我们可以先强调这个定义有多么模糊。人们一般认为猫就是两只猫交配后产生的东西，换句话说，猫的双亲都是猫。到这里为止还没有什么惊人的东西，甚至这可能还是非常显然的东西。但其中大有深意。
比如说现在有一只猫，它的双亲都是猫，它的双亲的双亲也是猫，它的双亲的双亲的双亲也是猫，以此类推。但是这种对时间的回溯并没有尽头!根据这种推理，我们必须回溯生命演化树，到达猫还没有出现的年代!的确，如果我们回溯到几亿年甚至几十亿年前，我们会到达一个没有哺乳动物、脊椎动物，甚至连真核生物都不存在的时代!猫的双亲都是猫，或者说猫就是两只猫交配后产生的东西，这些说法在逻辑上就有矛盾[10]。
我预计到有些读者会提出用DNA的概念来定义猫这个概念。然而这会导致几个问题。首先，我们还没有对猫的所有基因组进行测序，而我们可能永远无法对所有猫的基因组进行测序，因为未来的猫现在还没有出生!所以怎么定义一组与猫对应的DNA代码并不是一个显然的问题。其次，即使我们成功做到了这一点，还要对动物进行基因组测序才能断定它是不是猫，这个解决方案在现实中完全不可行。再次，人们也许会严肃地考虑，包含猫的全套DNA的一个单独的猫细胞，比如猫毛中的细胞，它算不算一只猫?最后也是最重要的一点就是，DNA这个概念在大约半个世纪之前完全不存在。往好了说，这意味着薛定谔、达尔文和亚里士多德在谈论猫的时候其实不知道自己在说什么。
我们必须理解这个显而易见的事实:“猫”这个词并没有一个令人满意的严谨的定义。如果我问你猫是什么，你无法给出一个普遍适用、毫无争议的定义。而这不是什么惊人的事情，毕竟你不是通过这种方式学到怎么认知并使用“猫”这个概念的!如果你大概知道猫是什么，那是因为你已经看过上千张甚至上百万张猫的图像，你也听说过猫，也读过有关猫的材料。你通过观察大量数据学到了猫是什么。但你从来没有看见过关于“猫是什么”的形式定义!实际上，在人类历史中，没有人知道怎么给出“猫是什么”的形式定义。
还有更神奇的。必定存在第一个想到“猫”这个概念的人。因为这个人是第一个想到这个概念的，不会有人教他这一点，所以这个人就自己发明了“猫”这个概念。为什么?他又是怎么做到的?人类引入的新概念从何而来?这种能力是人类智慧特有的吗?
我觉得谷歌人工智能的发现非常震撼人心，因为它给出了这些问题的答案。不，这并非人类智慧特有的能力，因为谷歌的神经网络同样做到了这一点。而它做到这一点是因为它希望分析、综合并解释数字化图像中不同像素颜色之间的相关性，它希望得到一个模型来描述自己看到的东西，而它找到的模型自然导致了“猫”这个概念的诞生!


正确问题的近似解答的价值远远大于错误问题的精确解答的价
值。----约翰·图基(1915—2000)
真理(......)实在过于复杂，除了近似以外都不可行。--约翰·冯·诺伊曼(1903—1957)

无穷计算在数学里遍地都是。这些计算中最著名的可能是周角率(它与历史上“喧宾夺主”的圆周率[2]的关系是)。在14世纪，伟大的印度智者摩陀婆就发现了一个惊人的等式
的交错和的8倍!;也就是说，周角率不过就是奇数倒数
我们可能会认为这种无穷计算毫无用处。但事实上，无穷计算在应用数学的简单模型中无处不在，流体力学方程(导数和积分都是无穷步的计算结果)就是一个例子。这是因为，即使无穷计算只代表了一种不可计算的理想情况，但它一般能指出进行近似计算的巧妙方法。比如摩陀婆的等式就能用于计算的近似值，也就是说，它可以用于计算半径已知的圆的周长。工程师在实践中使用的正是这类近似方法。
如果你在常用的计算器或者谷歌上查询的值，它们很有可能会撒一点小谎，只给你提供小数点后十几位数字的近似值。这个问题并非只在上发生。你的计算器对于那些二进制表示并非有限的数都处理得非常糟糕[3]，这样的数包括无理常数，例如和，还有某些有理数，比如1/3和0.2。此外，因为计算器只能处理近似值，所以它可能会得出一些违反常理的结果，比如说，或者在比大得多的时候会得到，即使本身大于0。
数学家经常会认为，在计算机上进行的计算只不过是数学理论的近似。所罗门诺夫妖的立场却完全相反。对于所罗门诺夫妖来说，实数之类的对象只不过是一些让我们可以构建算法并对其更好地进行思考的模型。特别是对于尝试追随所罗门诺夫妖步伐的计算机科学家来说，他们用于尝试衡量置信度的模型并不是由实数作为参数组成的模型，而是那些储存在计算机文件中的模型，其中用到的只是数学模型中实数的截断。正因如此，与数学的理想状态相反，人工神经网络的大小可以用字节来计数4。

渐近展开
如果说对这样的数值进行高精度计算毕竟获益有限，那么对物理学、生物学或者数学中出现的曲线、函数和行为进行近似就会产生大量的应用。通过素数定理对素数计数函数进行近似就是这样的情况。
更一般的框架中存在这样一项工具，它对于任意模型都可以计算出比模型本身简单得多的近似版本。这项工具就是渐近展开。一般来说，渐近展开可以让我们用一条直线来逼近圆上的一小段圆弧，或者用平面来逼近(相当圆润的)地球的一小块表面。用代数术语来说，这相当于用形如的仿射方程来逼近所谓的非线性方程。对于那些变化足够小的现象来说，这种近似完全是可以接受的。这就解释了为什么我们要在学校里花上这么长的时间来研究这些简单的方程，以及为什么这些方程在科学中无处不在5。

对物理学家来说不可或缺的渐近展开，最精彩的例子可能就是爱因斯坦的广义相对论方程。它的渐近展开可以导出牛顿力学!换个说法，就是牛顿的力学定律，特别是关于万有引力的定律，只不过是爱因斯坦方程的近似，在有限的场景里完全适用。这种有限的场景就是所谓的“弱引力”情况6。[4]
6更准确地说，这应该是时空曲率极小的情况。

实用主义的限制
然而在现实中，计算能力是有限的。我们之前也看到了，在地球上执行超过10^90个计算步骤是种不切实际的愿望。这就大大降低了我们用贝叶斯公式进行计算的能力。

作为比较，人类大脑包含大约10^15突触，这意味着对大脑的完整描述至少需要比特的信息。研究所有包含这么多字符的理论，至少需要2^{10^15}次计算!因此，即使能用到古戈尔个宇宙，对这些理论应用贝叶斯公式也完全是一种幻想。

图灵的机器学习
1950年，艾伦·图灵将这个关于算法不可避免的复杂度的论证过程漂亮地转移到了人工智能问题上。在一篇发表在期刊《心灵:心理学与哲学季度评论》(Mind,aQuarterlyReviewofPsychologyandPhilosophy)上，题为《计算机制与智能》(“ComputingMachineryandIntelligence”)的论文中，图灵首先提出了一个与算法执行所需时间没什么关系的问题:“机器能思考吗?”图灵尝试回答的就是这个问题。然而“思考”这个词的模糊性让他尝试将问题精确化。与其考虑这个问题，图灵提出了另一个问题，就是机器是否能像人那样行动。

更准确地说，图灵提出了以下的测试:要求一位人类A与另外两个实体X和Y进行书面通信，而这两个实体X和Y分别是人类和一台机器。如果人类A无法分辨X和Y之中哪一个实体是人类的话，那么机器就通过了测试。换句话说，如果机器知道怎么模仿人类，而任何人类都无法分辨出这台机器不是人类的话，那么机器就成功通过了测试。这就是图灵所谓的“模仿游戏”，我们今天将它称为图灵测试[5]。

在论文的第3节到第5节中，图灵重提了他1936年的工作，严格定义了机器到底是什么——这个概念推动了计算机的发明，接着就是数字时代的到来!然后在论文的第6节中，图灵否定了尝试证明机器不可能思考的9个经典论证。但更重要的是，在论文的第7节，图灵预料到了其测试中的难点及解决方法。即使仍然不能说计算机当时已经存在，但图灵不仅已经预料到了它们会在未来出现，而且还预料到了它们以后的性能将足以在模仿游戏中取胜:“(工程上的进步)似乎不可能不足够(使其通过测试)。”对于图灵来说，“这个问题主要是编程问题”。
特别是，图灵以惊人的远见推测出，能成功通过图灵测试的程序代码至少需要大概10^9个字符来描述。也就是说，用我们在第7章引入的术语来说，图灵自己猜测，图灵测试的所罗门诺夫复杂度应该以十亿字节来计量。为了得到这个估计值，图灵依靠的是他所知的唯一一台能通过图灵测试的机器。对，我说的就是人类大脑!毕竟还有什么能比人类更善于模仿人类呢?对图灵来说值得庆幸的是，当时的神经科学已经给出了人类大脑复杂度的估计值，提出人类的神经元之间大约有10^11到10^15个突触——现代的估计值处于10^14和5*10^14之间。
图灵提出，只有一小部分突触对于通过图灵测试来说是必不可少的，这就是需要10^15个字符这个数字的来历。

同样出于他无比的智慧，图灵注意到人类大脑能够通过图灵测试。于是他猜测，通过模仿人类大脑成功完成图灵测试，我们能更好地构建能够通过测试的机器。然而图灵也注意到，儿童教育同样是大脑最终发育过程中的重要一环。“儿童的大脑可能类似于我们在文具店买到的笔记本，”图灵这样写道，“没什么机关，但有许多张白纸8。”图灵从而提出，为了帮助机器填满它自己的白纸，可以让它从数据中学习。学习机器的概念，也就是能够学习的机器，就此诞生。

因此，学习机器这个想法能够让机器自己写出包含数十亿字符的程序——必要时可以写得更长。用更贴近算法的语言来说，这相当于肯定了机器学习最终可以让我们研究并探索那些描述长度超过十亿比特的算法。而更关键的是，引导这种探索的并非程序员的指尖，而是原始数据，就像儿童接收到的那样。
需要特别指出，图灵的这段论证反驳了众多知识分子的想法，其中还包含一些专家，他们经常这样说:“机器学习效果不错——但数学家不知道为什么。”这是2015年《连线》(Wired)杂志上一篇文章的标题。然而早在1950年，艾伦·图灵就预言了机器学习未来会取得成功，甚至明确指出它会在20世纪末出现!更厉害的是，关于机器学习会在众多任务中超越人类编写的程序这一情况，图灵已经强调了它的确切原因:只有机器学习才能探索那些长度超过十亿字节的算法空间[6]。
同样，某些专家经常指责那些庞大的神经网络无法被明确解释。然而，巨大而高效的神经网络无法用寥寥几个字符描述出来，至少无法以准确的方式描述出来，这并不令人意外。这是因为，如果神经网络能被简单描述的话，这个不太长的描述就会是简短的一种算法，它能够生成另一个算法(神经网络)，并借此解决图灵测试之类的问题。
这样的话，这个简短的算法就能解决图灵测试。然而，我们刚才推测，这项测试无法被简短的算法解决。

现在剩下的就是明确指出用什么方法来探索长度巨大的算法组成的空间。然而艾伦·图灵并没有指出要做到这一点可以使用的方法，他只是断言学习能力将会是关键。在21世纪初，人们提出了大量的方法。我们之前已经谈到过其中几种。不分先后的话，我们可以举出线性回归、逻辑斯谛回归、决策树、决策森林、支持向量机、神经网络、贝叶斯网络，还有马尔可夫随机场。我们会在本章和后面几章中再详细讨论最后这三个机器学习算法。

实用贝叶斯主义
为了解决像图灵测试那样拥有巨大的所罗门诺夫复杂度的问题，实用贝叶斯主义者只能放弃贝叶斯公式。他更应该利用那些所需时间没有超出现实可能性的方法。由此我们需要调整“有用”这个概念。对于纯粹贝叶斯主义者来说，置信度高的理论就是有用的;但对于实用贝叶斯主义者来说，置信度很高，但所需计算时间超出常理的理论实际上并没有用。因此，实用贝叶斯主义者将注意力转向了所需计算时间较短的理论之中更可信的那一部分。

这就让我们能够解释素数定理给出的近似以及牛顿运动定律在实用中的置信度了。对于纯粹贝叶斯主义者来说，这两种描述都没有任何用处，因为两者分别被素数的精确计算以及爱因斯坦的广义相对论超越。然而实用贝叶斯主义者则乐于拥抱这些近似方法，因为它们对计算的需求明显小得多。计算的近似值需要的时间是的对数，而牛顿理论中的向量和微分计算与爱因斯坦理论中的张量计算相比远远更快。只要在我们身处的场景中，这两种近似精确到足以被接受(分别是值很大的情况和弱引力的情况)，那么对于实用贝叶斯主义者来说，这两种近似就比精确的版本更有用!
这种与计算时间有关的“有用”概念同样可以作为对神经网络惊人的成功的首要解释。这是因为，神经网络，或者至少是所谓的前馈神经网络，与那些更精细的算法相反，需要的计算时间不长，而且必定有所限制(图14.1)。实际上，如果我们考虑前馈神经网络的一个足够广泛的定义，那么这些网络组成的集合恰好就是快速并行算法的集合。因此，在前馈神经网络上进行机器学习，其实就相当于利用快速算法尽可能好地对数据进行解释(如果可能的话，也要利用合适的贝叶斯先验置信度)。所以，这就是迈向实用贝叶斯主义的第一步。
<!-- /var/folders/gl/1x3lhvs57l72x96lt82xqx_r0000gp/T/TemporaryItems/NSIRD_screencaptureui_K6o907/截屏2024-09-0400.00.07.png -->

图14.1神经网络就是神经元之间一系列通信以及神经元本身进行的基本计算的总和。如果这种通信是无环的，我们就说这是个前馈神经网络，从直觉上来说，也就是通信只会朝一个方向进行，从来不会形成一个闭环



思



人工智能(ArtificialIntelligence,AI)始于一个简单的想法。
计算机被发明之后，人们很快认识到它的能力远不限于数值计算，可以通过对信息和指令编程完成很多以往只有人脑才能完成的任务。


强调像人一样
强调合理性
思考

像人一样地思考

合理地思考

行动
像人一样地行动
合理地行动


||强调像人一样|强调合理性|
|:----:|:----:|:----:|
|思考|像人一样地思考|合理地思考|
|行动|像人一样地思考|合理地行动|


###人工智能：广跨各学科，综合文理工

人工智能：对人的意识、思维过程进行模拟的一门学科。

###发展历程：AI的三次热潮和相应的研究学科突破

人工智能发展如同过山车，忽上忽下，从创立以来不断分化

AI数理基础切换:从基于逻辑的表达与推理到基于统计的学习与计算

###三次热潮的总结

第一阶段(前30年)：以数理逻辑的表达与推理为主。
杰出代表，如JohnMcCarthy、MarvinMinsky、HerbertSimmon：懂很多认知科学的东西，有很强的全局观念，拿过图灵奖和很多大奖。但工具基本都是基于数理逻辑和推理。
逻辑的东西发展得干净、漂亮。但符号的知识表达不落地；


 第二阶段(近30年)：以概率统计的建模、学习和计算为主。
视觉、自然语言、认知、机器学习、机器人五大领域在1990年中期都找到了概率统计这个新“机制”：统计建模、机器学习、随机算法……


###人工智能的主要学派
|符号主义学派Symbolicism|联结主义学派Connectionism|行为主义学派Actionism|
|:----:|:----:|:----:|
|起源于2300年前的逻辑主义、心理学派和计算机学派|起源于上世纪40年代仿生学派或生理学派|起源于上世纪40年代，进化主义、控制论|
|认知即计算：用符号、规则和逻辑来表征指示和进行逻辑推理|认知即网络：用概率矩阵和加权神经元进行模式识别、分类|认知即行动：基于“感知-行动”生成变化，为特定目标获取其中最优解|
|基本算法：规则和决策树|基本算法：深度学习神经网络|基本算法：遗传算法、蚁群算法、强化学习|

###三大学派的方法论

|符号主义学派Symbolicism|联结主义学派Connectionism|行为主义学派Actionism|
|:----:|:----:|:----:|
|知识自上而下地设计|知识自下而上地涌现|知识自下而上地涌现|
|经由专家自上而下基于逻辑设计、推理来实现|通过模拟大脑的结构(神经网络)来实现知识|从简单生物体与环境互动的模式中寻找知识







参考文献

[1]https://blog.csdn.net/2401_84033492/article/details/137194675

[2]https://blog.csdn.net/xierhacker/article/details/53431207

[3]https://blog.csdn.net/u014313009/article/details/51039334

[4]https://baike.baidu.com/item/%E4%BA%BA%E5%B7%A5%E7%A5%9E%E7%

BB%8F%E7%BD%91%E7%BB%9C/382460#4

[5]https://zh.wikipedia.org/wiki/**最邻近搜索**

